% This file is part of the CTAN package named isIntDim.
% 
%   info.tex: instructions for two macro packages: ifisinteger.tex
%             and ifisdimension.tex
%
%   Copyright (C) 2024  Udo Wermuth (author)
%
%   This program is free software: you can redistribute it and/or modify
%   it under the terms of the GNU General Public License as published by
%   the Free Software Foundation, either version 3 of the License, or
%   (at your option) any later version.
%
%   This program is distributed in the hope that it will be useful,
%   but WITHOUT ANY WARRANTY; without even the implied warranty of
%   MERCHANTABILITY or FITNESS FOR A PARTICULAR PURPOSE.  See the
%   GNU General Public License for more details.
%
%   You should have received a copy of the GNU General Public License
%   along with this program.  If not, see <http://www.gnu.org/licenses/>.

% %%%
% %%% verbatim macros (from manmac.tex)
% %%%
\newskip\ttglue {\tt\global\ttglue=0.5em plus 0.25em minus 0.15em }
\def\ttverbatim{\begingroup \frenchspacing
  \catcode`\\=12 \catcode`\{=12 \catcode`\}=12 \catcode`\$=12 \catcode`\&=12
  \catcode`\#=12 \catcode`\%=12 \catcode`\~=12 \catcode`\_=12 \catcode`\^=12
  \obeyspaces \obeylines \tt}
\def\verbatimspace{\ifvmode\indent\fi\space}
{\obeyspaces \gdef\makespaceverbspace{\def {\verbatimspace}}}
\outer\def\verbatim{$$\ifdim\parskip>0pt
    \abovedisplayskip=\parskip \abovedisplayshortskip=\parskip
    \belowdisplayskip=\parskip \belowdisplayshortskip=\parskip
  \else
    \abovedisplayskip=3pt \abovedisplayshortskip=3pt
    \belowdisplayskip=3pt \belowdisplayshortskip=3pt
  \fi
  \let\par=\endgraf \ttverbatim \makespaceverbspace \parskip=0pt
  \catcode`\�=0 \advance\leftskip by 10pt \ttfinish}
{\catcode`\�=0 �catcode`�\=12 % � is temporary escape character
  �obeylines % end of line is active
  �gdef�ttfinish#1^^M#2\endverbatim{�vbox{#2}�endgroup$$}}
\catcode`\|=\active
{\obeylines \gdef|{\ttverbatim \spaceskip\ttglue \let^^M=\  \let|=\endgroup}}

\def\noitem{\item{\phantom{0.}}}

% %%%%%%%%%%%%%%%%%%%%%%%%%%%%%%%%%%%%%%%%%%%%%%%%%%%%%%%%%%%%%%%%%%%%%%%%%%%%
\font\titlefont=cmssdc10 at 36pt
\font\subtitlefont=cmssdc10 at 17pt
%
\centerline{\titlefont ifis-macros}
\bigskip
\centerline{\subtitlefont Version 1.0, 7.04.2024}
\bigskip
\centerline{Macros for plain \TeX}
\medskip
% %%%%%%%%%%%%%%%%%%%%%%%%%%%%%%%%%%%%%%%%%%%%%%%%%%%%%%%%%%%%%%%%%%%%%%%%%%%%

\noindent
There are two main macros in the files {\tt ifisinteger.tex} and {\tt
ifisdimension.tex}. The macro |\ifisint| tests if a given input string
represents a number for \TeX. The macro |\ifisdim| does this for
dimensions.

Both macros generate errors but hide them from the terminal as they
work in |\batchmode|. There is one configuration parameter:
|\IIcurrentmode|. The default is |\errorstopmode|. Change this if you
call the macros in a different interaction mode so that they can
return to this mode.

\bigskip

\beginsection 0.\ Installation

To use the macro |\ifisint| load via |\input ifisinteger.tex| the file
that contains the code. For |\ifisdim| use |\input ifisdimension.tex|.

\bigskip

\beginsection 1.\ File {\tt ifisinteger.tex}

The main macro is called |\ifisint| and must be used like an
|\if|-conditional except that its argument is delimited by |\Boolend|:
|\ifisint <argument>\Boolend <true branch>\else <false branch>\fi|.

\medskip

The implemented algorithm has four steps:
\smallskip
\item{1.}1)~Remove signs with or without braces; add sentinel |W|.
\noitem  2)~Test that the input isn't now |"W|, etc.; otherwise
         return false.
\item{2.}Create a canonical form with a leading zero.
\item{3.}1)~Assign the input to a |\count| register inside an
         |\hbox|.
\noitem  2)~Test that the box width is the width of the sentinel.
\noitem  3)~Otherwise return false.
\item{4.}1)~Return true if the number isn't \TeX's maximum.
\noitem  2)~Otherwise test if the canonical form is \TeX's maximum.
         If yes, return true.
\noitem  3)~Otherwise return false.
\medskip

For more details see my article in TUGboat {\bf45}:1 (2024), 106--109.

\bigskip

\beginsection 2.\ File {\tt ifisdimension.tex}

The main macro is called |\ifisdim| and must be used like an
|\if|-conditional except that its argument is delimited by |\Boolend|:
|\ifisdim <argument>\Boolend <true branch>\else <false branch>\fi|.

\medskip

The implemented algorithm has four steps:
\smallskip
\item{1.}1)~Remove signs with or without braces; add sentinel |mm|.
\noitem  2)~Exclude trivial non-numerics as done in |\ifisint|.
\noitem  3)~Otherwise return false.
\item{2.}1)~Get the integer part.
\noitem  2)~Get fraction and the unit.
\noitem  3)~Get the width of the unit.
\item{3.}1)~Assign the input to a |\dimen| register inside an
         |\hbox|.
\noitem  2)~Test that the box width is the width of the sentinel.
\noitem  3)~Otherwise return false.
\item{4.}1)~Return true if the dimension isn't \TeX's |\maxdimen|.
\noitem  2)~Otherwise test if the coerced sum of the integer part
         and the fraction expressed in the unit |sp| is |\maxdimen|.
\noitem  3)~If no, return false.
\noitem  4)~Otherwise return true.
\medskip

For more details see my article in TUGboat {\bf45}:1 (2024), 109--112.

\bye
